\subsection{北美 AI 硬件与芯片安全研究现状}

北美 AI 硬件与芯片安全研究现状的核心特点可以概括为“漏洞发现快、产业转化快、云端闭环强”。学术界长期把微架构侧信道、GPU 内存隔离、多 GPU 互连、硬件木马和可信执行环境纳入系统安全问题来研究；产业界则由 NVIDIA、AMD、Intel、Apple、云厂商和安全响应机构共同推动，将真实 GPU、CVM、远程证明、固件签名、驱动修复和端云协同 AI 平台纳入产品化闭环。

\begin{table*}[htbp]
\centering
\scriptsize
\renewcommand{\arraystretch}{1.25}
\begin{tabularx}{\textwidth}{p{0.9cm}p{1.75cm}p{1.9cm}X X X X}
\toprule
核心维度 & 细分落地 & 核心主体 & 代表性成果 & 技术方案 & 核心价值 & 不足 \\
\midrule
学术侧 & GPU 渲染、浏览器与微架构侧信道 & UC 系、University of Texas、University of Chicago 等；GPU.zip、Hot Pixels、WebGPU-SPY 团队 &
1. GPU.zip 证明图形数据压缩可形成跨源像素泄露侧信道 \cite{gpuzip2024}。\newline
2. Hot Pixels 展示频率、功耗、温度等混合信号可支持像素窃取和历史嗅探 \cite{hotpixels2023}。\newline
3. WebGPU-SPY 将 GPU cache 攻击引入 WebGPU 沙箱 \cite{webgpuispy2024}。 &
1. 利用 GPU 压缩、渲染延迟、像素布局差异构造跨源观测。\newline
2. 通过 DVFS、温度、功耗传感器推断图形内容。\newline
3. 借助 WebGPU 高并行计时与 cache 争用建立浏览器侧信道。 &
1. 说明浏览器 Agent 的底层风险包含 GPU 渲染路径。\newline
2. 推动浏览器、驱动和 GPU API 重新评估隔离粒度。\newline
3. 为端侧多模态 Agent 和浏览器 Agent 攻击面建模提供依据。 &
1. 攻击效果依赖具体 GPU、驱动、浏览器和页面布局。\newline
2. 缓解措施可能牺牲图形性能或 WebGPU 可用性。\newline
3. 与企业浏览器和大模型容器环境还需更多验证。 \\
\midrule
学术侧 & GPU 显存、本地内存与多租户 AI 算力隔离 & Trail of Bits、UC Santa Cruz、CMU CERT/CC；多 GPU 侧信道研究团队 &
1. LeftoverLocals 证明 AMD、Apple、Qualcomm 等 GPU local memory 清理不足可泄露 LLM 响应 \cite{leftoverlocals2024,certleftoverlocals2024}。\newline
2. Spy in the GPU-box 证明 NVIDIA DGX 多 GPU 系统的远端 GPU L2 cache 争用可形成隐蔽信道和侧信道 \cite{spygpubox2022}。\newline
3. NVBleed/Beyond the Bridge 将 NVLink/PCIe 互连纳入云端多 GPU 泄露面分析 \cite{nvbleed2025,beyondbridge2024}。 &
1. 读取 GPU kernel 结束后残留的 local memory。\newline
2. 逆向多 GPU cache/互连层次，通过 contention、prime-probe、性能计数器观测远端负载。\newline
3. 在多租户或虚拟化场景中测试互连拥塞、缓存命中和应用指纹泄露。 &
1. 将 prompt、响应、中间张量和模型片段纳入 GPU 隔离威胁模型。\newline
2. 说明共享 GPU 不是普通共享 CPU 资源。\newline
3. 为云 GPU 分配、容器隔离、MIG/vGPU 策略和清零机制提供依据。 &
1. 复现依赖特定设备、API、驱动版本和调度条件。\newline
2. 清零、禁用并发、限制计数器会带来吞吐损失。\newline
3. 云厂商真实多租户策略和硬件细节不完全公开。 \\
\midrule
学术侧 & 硬件木马、供应链可信与安全 EDA & University of Florida、NYU、UC 系、DARPA AISS 相关高校/企业团队 &
1. 硬件木马分类与检测研究将威胁扩展到 IP 复用、设计外包和供应链可信 \cite{tehranipoor2010trojan}。\newline
2. 信息流安全验证尝试在 RTL/IP 层检测恶意逻辑或违反安全策略的设计 \cite{nahiyan2018ifs}。\newline
3. DARPA AISS 将侧信道、逆向工程、供应链攻击和恶意硬件纳入自动化安全芯片设计流程 \cite{darpaAISS}。 &
1. 建立硬件木马 taxonomy、触发机制、载荷类型和检测方法。\newline
2. 在设计阶段使用信息流、形式化验证、测试向量和侧信道分析识别异常。\newline
3. 将安全分区、IP 完整性和安全成本权衡融入 SoC 自动化流程。 &
1. 回答芯片、IP、固件和供应链是否可信的底层问题。\newline
2. 将安全从部署后补丁前移到芯片设计与验证阶段。\newline
3. 为 AI 加速器、NPU、DPU 和定制 SoC 的安全认证提供基础方法。 &
1. 与具体 AI 加速器和大模型负载结合仍不充分。\newline
2. 面对大规模 SoC 时成本高、覆盖有限。\newline
3. 供应链可信依赖多方协同，落地复杂。 \\
\midrule
产业侧 & 云端机密 AI 加速与 CPU-GPU TEE 协同 & NVIDIA、AMD、Intel、Google Cloud、Microsoft Azure 等 &
1. NVIDIA H100 将 GPU 纳入机密计算边界，提供 RoT、测量启动、SPDM、设备证明、CC-On 和加密 CPU-GPU 传输 \cite{nvidiah100cc2023}。\newline
2. AMD SEV-SNP 和 Intel TDX 为机密 VM 提供内存加密、完整性保护和远程证明 \cite{amdsev2026,inteltdx2026}。\newline
3. NVIDIA Secure AI/NRAS 将 GPU 证明、密钥发布、镜像治理和企业 AI 工作流连接起来 \cite{nvidiasecureai2025}。 &
1. 在 GPU 内建立硬件 RoT、固件签名、测量启动和设备证书链。\newline
2. 用 CVM 保护主机侧内存与驱动路径，用加密 bounce buffer/命令签名保护 PCIe 传输。\newline
3. 远程验证 GPU、驱动、固件、镜像和运行模式后再释放数据或模型。 &
1. 保护 AI 训练/推理中的数据、模型权重和企业上下文。\newline
2. 支撑敏感 Agent 在公有云/混合云部署。\newline
3. 把硬件安全能力转化为企业 AI 平台能力。 &
1. TCB 跨 CPU、GPU、固件、驱动、证明服务和云编排。\newline
2. PCIe 加密、数据搬运和证明流程可能带来性能与运维成本。\newline
3. 关键实现依赖闭源固件、证书服务和云平台策略。 \\
\midrule
产业侧 & GPU 漏洞响应、驱动修复与硬件安全治理 & AMD、Apple、Qualcomm、Google、Khronos、CERT/CC、Trail of Bits &
1. CERT/CC VU\#446598 对 LeftoverLocals 进行协调披露，确认 AMD、Apple、Qualcomm 等 GPGPU 平台存在隔离不足风险 \cite{certleftoverlocals2024}。\newline
2. AMD-SB-6010 将问题归入 GPU Memory Leaks，并给出 CVE-2023-4969、驱动/固件修复和可选新运行模式 \cite{amdleftoverlocals2024}。\newline
3. Apple、Google 等围绕受影响设备、ChromeOS/芯片代际和软件更新路径进行修复或缓解 \cite{certleftoverlocals2024}。 &
1. 通过驱动、固件、运行模式限制并发 GPU 进程并清理寄存器/local memory。\newline
2. 以 CVE、供应商公告、标准组织和协调披露机制推动修复。\newline
3. 对 OpenCL、Metal、Vulkan、ROCm、系统驱动和设备固件进行版本化治理。 &
1. 把学术/安全公司发现快速转化为厂商修复路线。\newline
2. 强化 GPU local memory、寄存器、kernel 边界的安全基线。\newline
3. 促使 GPU 厂商把多租户 AI、容器和 LLM 场景纳入安全公告。 &
1. 修复链条涉及芯片商、设备商、OS、驱动、云平台和用户升级。\newline
2. 可选隔离模式可能降低并发能力和利用率。\newline
3. 老旧设备、移动 SoC 和闭源驱动的长期修复可达性不稳定。 \\
\midrule
产业侧 & 端云协同私有 AI 与可验证云推理节点 & Apple Private Cloud Compute；Apple Silicon、Secure Enclave、Apple Intelligence 生态 &
1. Apple PCC 将 Apple Silicon、Secure Enclave、Secure Boot、代码签名和受限 OS 引入云端 AI 推理节点 \cite{applepcc2024}。\newline
2. PCC 提出无远程 shell、无通用日志、无特权运行时访问和请求处理后不保留用户数据 \cite{applepcc2024}。\newline
3. PCC 通过透明日志、公开生产软件镜像和设备侧证明，让用户设备只向可验证节点发送请求 \cite{applepcc2024}。 &
1. 端侧设备先验证 PCC 节点证书、系统测量值和透明日志，再加密发送请求。\newline
2. 通过 Secure Boot、代码签名、trust cache、SEP 密钥保护和内存安全语言缩小攻击面。\newline
3. 通过无特权接口、受限观测、请求后清除和 target diffusion 降低运营方与定向攻击风险。 &
1. 将端侧隐私模型延伸到云端大模型推理。\newline
2. 从“云服务承诺不看数据”转向“硬件、镜像和证明链共同约束服务提供者”。\newline
3. 对移动端 Agent、个人智能助理和隐私敏感推理具有示范意义。 &
1. 强依赖 Apple 封闭硬件、OS 和服务生态。\newline
2. 关键硬件与部分实现仍由厂商控制。\newline
3. 主要解决云端推理隐私，不直接处理提示注入、工具越权和长期日志治理。 \\
\bottomrule
\end{tabularx}
\caption{北美 AI 硬件与芯片安全研究现状：以方向为主轴、主体为支撑的结构化总结}
\label{tab:north-america-ai-hardware-security}
\end{table*}

学术侧对北美 AI 硬件安全的推进，突出体现在对“真实 GPU 行为”的攻击面重估。GPU.zip、Hot Pixels 和 WebGPU-SPY 共同说明，现代 GPU 的图形压缩、渲染流水线、DVFS、传感器、cache 和浏览器 GPU API 都可能变成可观测信道 \cite{gpuzip2024,hotpixels2023,webgpuispy2024}。这对浏览器 Agent 尤其关键：攻击者未必需要突破应用层权限，只要能够诱导页面、渲染路径或 GPU API 进入特定状态，就可能间接观察跨源视觉信息、页面指纹或用户行为。

GPU 内存隔离和多租户 AI 算力隔离已经成为北美研究中最贴近大模型部署的方向。LeftoverLocals 把 GPU local memory 清理不足与 LLM 响应泄露直接联系起来，说明 prompt、响应、中间张量和模型片段都可能在 GPU kernel 与进程边界之间残留 \cite{leftoverlocals2024,certleftoverlocals2024}。Spy in the GPU-box、Beyond the Bridge 和 NVBleed 则把问题推进到多 GPU 系统，证明 NVLink、PCIe、远端 cache 和互连拥塞并不天然安全 \cite{spygpubox2022,beyondbridge2024,nvbleed2025}。

北美在更传统的硬件安全方向上积累深厚，硬件木马、供应链可信和安全 EDA 构成 AI 芯片底座安全的前置问题。硬件木马分类与检测工作将风险从制造阶段扩展到 IP 复用、设计外包和测试不足等环节 \cite{tehranipoor2010trojan}；信息流安全验证等研究进一步尝试在设计阶段识别恶意逻辑和安全策略违背 \cite{nahiyan2018ifs}。DARPA AISS 则把侧信道、逆向工程、供应链攻击和恶意硬件纳入自动化安全芯片设计流程，体现出“安全机制前移到 SoC 设计阶段”的政策与工程导向 \cite{darpaAISS}。

产业侧最强的主线是云端机密 AI 加速。NVIDIA H100 将 GPU 纳入机密计算边界，通过硬件根信任、测量启动、SPDM、设备证明、CC-On 模式、硬件防火墙和加密 CPU-GPU 数据通道，保护 AI 数据和模型在使用中的机密性 \cite{nvidiah100cc2023}。AMD SEV-SNP 和 Intel TDX 则为主机侧机密 VM 提供内存加密、完整性保护和远程证明，使 CPU TEE 与 GPU TEE 可以共同支撑云端敏感 AI 工作负载 \cite{amdsev2026,inteltdx2026}。

产业侧对漏洞响应的速度也体现出北美生态的闭环能力。LeftoverLocals 从研究发现、CERT/CC 协调披露到 AMD-SB-6010 厂商公告，形成了从 PoC 到 CVE、驱动/固件修复和可选运行模式的转化链条 \cite{certleftoverlocals2024,amdleftoverlocals2024}。这种机制能够较快把安全研究转化为产品修复，但 GPU 生态链条长，涉及芯片商、设备商、OS、驱动、标准组织、云平台和终端用户，老旧设备和移动 SoC 的修复可达性并不稳定。

Apple Private Cloud Compute 代表了北美端云协同 AI 安全的另一条路线。PCC 将 Apple Silicon、Secure Enclave、Secure Boot、代码签名、无远程 shell、受限日志、透明日志和设备侧证明带入云端 AI 推理节点，使用户设备只向可验证的软件镜像和硬件节点发送请求 \cite{applepcc2024}。它的意义在于从“服务商承诺不访问数据”转向“硬件、系统镜像、证明链和透明日志共同约束服务提供者”。但 PCC 主要解决云端推理隐私和运营方访问控制，不能替代应用层 Agent 安全。

总体看，北美路线对 Agent 安全的启示是：Agent 运行时应放在可证明、可度量、可审计的可信算力底座之上，但不能把硬件可信等同于 Agent 安全完备。机密 GPU、CVM、远程证明和硬件根信任可以降低恶意租户、云运营方和供应链攻击者的能力；浏览器隔离、显存清零、多 GPU 互连治理和端云透明证明可以减少底层泄露面；而提示注入、工具调用越权、权限扩散和日志治理仍需要与硬件底座协同设计。
